% SPDX-License-Identifier: Apache-2.0
\section{Late Binding, Tags and Configuration}
\label{sec:e-binding}

The merged language has a whole facet for deciding what is bound to what.
It holds \code{config}, \code{bind}, \code{domain}, \code{blackbox} and
\code{foreign}, and it exists so that changing a target edits one file and
leaves a verified design alone.

Embedded in Rust, the facet keeps its job and loses its syntax.

\subsection{A config is a struct}

A config gathers every choice a build makes: which implementation fills
each socket, what each generic parameter is, and what the physical numbers
are. A trait with associated types and associated constants holds all
three.

\begin{lstlisting}[caption={A config is a struct. Associated types hold the sockets and the domains, associated constants the parameters and the physical numbers.},label={lst:config}]
pub trait Config {
    type Mac: MacOp;          // a socket
    type Clock: Tag;          // a domain
    const TAPS: usize;        // a generic parameter
    const CLK_HZ: u64;        // a physical number
}

pub struct Fpga;
impl Config for Fpga {
    type Mac = DspSliceMac;
    type Clock = Elastic;
    const TAPS: usize = 16;
    const CLK_HZ: u64 = 100_000_000;
}
\end{lstlisting}

The design is generic over the config rather than over one parameter per
choice, so adding a knob does not change every unit's signature.

\subsection{The path in a binding disappears}

\code{bind top.filter.mac => DspSliceMac} names a path through the
instance hierarchy. That path is the part most likely to go stale, because
it breaks whenever the hierarchy is rearranged.

A nested unit is generic over the same config, so the choice reaches it
without anyone naming a path.

\begin{lstlisting}[caption={The path in a binding disappears. A nested unit is generic over the same config, so rearranging the hierarchy cannot invalidate anything.},label={lst:nesting}]
pub struct Filter<C: Config> { pub mac: C::Mac, pub acc: u64 }
pub struct Top<C: Config>    { pub filter: Filter<C> }

// The configuration facet, entire.
pub type FpgaBuild = Top<Fpga>;
pub type AsicBuild = Top<Asic>;
\end{lstlisting}

Rearranging the hierarchy cannot invalidate a binding, because there is no
binding to invalidate. The compiler checks that each implementation
satisfies its contract, which is what the specification said \code{bind}
had to do.

\subsection{The numbers are compile-time}

That claim is easy to make and easy to get wrong, so it was compiled. An
array length is a context that accepts nothing but a true constant:

\begin{lstlisting}[caption={Proof that the config's numbers are compile-time: an array length accepts nothing else.},label={lst:const}]
pub const FPGA_TAPS: usize = <Fpga as Config>::TAPS;
pub const TAP_BUFFER: [u64; FPGA_TAPS] = [0; FPGA_TAPS];
\end{lstlisting}

\subsection{Tags are types, and not for the reason expected}

A tag names a synchronisation domain and its policy. The natural spelling
is an attribute, \code{\#[tag(MemFetch, handshake, capacity = 8)]}, and it
is the spelling the merged specification chose.

It does not work on a block. An attribute on an expression is an unstable
feature, and the probe reports \code{error[E0658]}. A tag is therefore a
type with associated constants, and a unit is generic over it.

\begin{lstlisting}[caption={A tag is a type with associated constants, because an attribute on a block is not stable.},label={lst:tag}]
pub struct MemFetch;
impl Tag for MemFetch {
    const HANDSHAKE: bool = true;
    const CAPACITY: usize = 8;
}

pub struct MacUnit<T: Tag> { /* ... */ }
\end{lstlisting}

That was going to be the design regardless, because a type parameter is
what lets a config choose the domain. The probe turned a preference into a
reason.
